\documentclass[12pt,a4paper]{article}
\usepackage[utf8]{inputenc}
\usepackage[brazilian]{babel}
\addto\captionsbrazilian{
  \renewcommand{\lstlistingname}{Código}
  \renewcommand{\lstlistlistingname}{Lista de Códigos}
}
\usepackage[T1]{fontenc}
\usepackage{geometry}
\geometry{left=2.5cm,right=2.5cm,top=2.5cm,bottom=2.5cm}
\usepackage{listings}
\usepackage{xcolor}
\usepackage{graphicx}
\usepackage{hyperref}
\usepackage{float}
\usepackage{booktabs}
\usepackage{longtable}
\usepackage{amssymb}
\usepackage{amsmath}
\usepackage{microtype}
\usepackage{abstract}
\setlength{\absleftindent}{0pt}
\setlength{\absrightindent}{0pt}

\definecolor{codegreen}{rgb}{0,0.6,0}
\definecolor{codegray}{rgb}{0.5,0.5,0.5}
\definecolor{codepurple}{rgb}{0.58,0,0.82}
\definecolor{backcolour}{rgb}{0.95,0.95,0.92}

\lstdefinestyle{mystyle}{
    backgroundcolor=\color{backcolour},
    commentstyle=\color{codegreen},
    keywordstyle=\color{magenta},
    numberstyle=\tiny\color{codegray},
    stringstyle=\color{codepurple},
    basicstyle=\ttfamily\footnotesize,
    breakatwhitespace=false,
    breaklines=true,
    captionpos=b,
    keepspaces=true,
    numbers=left,
    numbersep=5pt,
    showspaces=false,
    showstringspaces=false,
    showtabs=false,
    tabsize=2
}
\lstset{style=mystyle}

\title{Compiladores -- Trabalho Prático (Parte 2) \\ Analisador Sintático Descendente Preditivo}
\author{Alunos: Luan Marcelino de Souza e Vitor Mendes de Oliveira Abreu \\ Prof.ª Kércia Aline Marques Ferreira \\ CEFET-MG}
\date{\today}

\begin{document}

\shorthandoff{"}

\maketitle

\begin{abstract}
Este relatório descreve a segunda etapa do trabalho prático da disciplina de Compiladores: o desenvolvimento de um analisador sintático preditivo descendente recursivo implementado manualmente em Java. A gramática original foi modificada para adequar-se às restrições de uma gramática $LL(1)$ com eliminação de recursão à esquerda e fatoração à esquerda. São apresentados os conjuntos $FIRST$ e $FOLLOW$, a arquitetura do software organizada nos pacotes \texttt{lexico}, \texttt{sintatico} e \texttt{tabela}, e os resultados da análise sobre os casos de teste.
\end{abstract}

\section{Introdução}

O analisador sintático é o componente do front-end de um compilador responsável por verificar se a sequência de tokens gerada pelo analisador léxico está em conformidade com a sintaxe formal da linguagem.

Nesta etapa, estendendo a infraestrutura da Etapa 1, implementou-se manualmente um parser descendente preditivo por recursão mútua ($LL(1)$), no qual cada símbolo não-terminal da linguagem torna-se um método Java. O fluxo de análise é determinístico: a escolha entre produções alternativas é feita unicamente com base no token corrente (lookahead), sem necessidade de retrocesso.

O código-fonte está organizado em três pacotes:
\begin{itemize}
    \item \texttt{lexico} --- classes do analisador léxico (\texttt{Lexer}, \texttt{Token}, \texttt{Tag}, \texttt{Word}, \texttt{Num}, \texttt{Literal})
    \item \texttt{sintatico} --- analisador sintático (\texttt{Parser})
    \item \texttt{tabela} --- tabela de símbolos (\texttt{Env}, \texttt{Id})
\end{itemize}

\section{Como usar o compilador (Fase Sintática)}

\subsection{Compilação}

\begin{lstlisting}[language=bash, caption=Compilação do analisador]
javac -d out compilador\src\Main.java compilador\src\lexico\*.java compilador\src\tabela\*.java compilador\src\sintatico\*.java
\end{lstlisting}

\subsection{Execução}

\begin{lstlisting}[language=bash, caption=Execução do parser]
java -cp out Main compilador\Teste7.txt
\end{lstlisting}

O programa exibe \texttt{"Programa sintaticamente correto."} em caso de sucesso, ou aborta ao encontrar o primeiro erro sintático, indicando a linha e o token inesperado.

\allowdisplaybreaks[4]

\section{Gramática da Linguagem}

A gramática da linguagem é apresentada a seguir, dividida em produções sintáticas (analisador sintático) e regras de formação de tokens (analisador léxico).

\subsection{Produções Sintáticas}

\begin{align*}
\hspace*{2em}1.&\quad program &&::= \text{class } identifier \; \text{\{} \; {[decl\!-\!list]} \; body \; \text{\}} \\
2.&\quad decl\!-\!list &&::= decl \; \text{;} \; \{ \; decl \; \text{;} \; \} \\
3.&\quad decl &&::= type \; ident\!-\!list \\
4.&\quad ident\!-\!list &&::= identifier \; \{ \; \text{,} \; identifier \; \} \\
5.&\quad type &&::= \text{int} \\
6.&\quad &&\mid \; \text{string} \\
7.&\quad &&\mid \; \text{float} \\
8.&\quad body &&::= \text{\{} \; stmt\!-\!list \; \text{\}} \\
9.&\quad stmt\!-\!list &&::= stmt \; \text{;} \; \{ \; stmt \; \text{;} \; \} \\
10.&\quad stmt &&::= assign\!-\!stmt \\
11.&\quad &&\mid \; if\!-\!stmt \\
12.&\quad &&\mid \; do\!-\!stmt \\
13.&\quad &&\mid \; repeat\!-\!stmt \\
14.&\quad &&\mid \; read\!-\!stmt \\
15.&\quad &&\mid \; write\!-\!stmt \\
16.&\quad assign\!-\!stmt &&::= identifier \; \text{:=} \; simple\!-\!expr \\
17.&\quad if\!-\!stmt &&::= \text{if } \text{(} \; condition \; \text{)} \; \text{\{} \; stmt\!-\!list \; \text{\}} \; if\!-\!stmt' \\
18.&\quad if\!-\!stmt' &&::= \text{else } \text{\{} \; stmt\!-\!list \; \text{\}} \\
19.&\quad &&\mid \; \lambda \\
20.&\quad do\!-\!stmt &&::= \text{do } \text{\{} \; stmt\!-\!list \; \text{\}} \; do\!-\!suffix \\
21.&\quad do\!-\!suffix &&::= \text{while } \text{(} \; condition \; \text{)} \\
22.&\quad repeat\!-\!stmt &&::= \text{repeat } \text{\{} \; stmt\!-\!list \; \text{\}} \; stmt\!-\!suffix \\
23.&\quad stmt\!-\!suffix &&::= \text{until } \text{(} \; condition \; \text{)} \\
24.&\quad read\!-\!stmt &&::= \text{read } \text{(} \; identifier \; \text{)} \\
25.&\quad write\!-\!stmt &&::= \text{write } \text{(} \; writable \; \text{)} \\
26.&\quad writable &&::= simple\!-\!expr \\
27.&\quad condition &&::= expression \\
28.&\quad expression &&::= simple\!-\!expr \; expression' \\
29.&\quad expression' &&::= relop \; simple\!-\!expr \\
30.&\quad &&\mid \; \lambda \\
31.&\quad simple\!-\!expr &&::= term \; simple\!-\!expr' \\
32.&\quad simple\!-\!expr' &&::= addop \; term \; simple\!-\!expr' \\
33.&\quad &&\mid \; \lambda \\
34.&\quad term &&::= factor\!-\!a \; term' \\
35.&\quad term' &&::= mulop \; factor\!-\!a \; term' \\
36.&\quad &&\mid \; \lambda \\
37.&\quad factor\!-\!a &&::= factor \\
38.&\quad &&\mid \; \text{not } factor \\
39.&\quad &&\mid \; \text{-} factor \\
40.&\quad factor &&::= identifier \\
41.&\quad &&\mid \; constant \\
42.&\quad &&\mid \; \text{(} \; expression \; \text{)} \\
43.&\quad relop &&::= \text{>} \\
44.&\quad &&\mid \; \text{>=} \\
45.&\quad &&\mid \; \text{<} \\
46.&\quad &&\mid \; \text{<=} \\
47.&\quad &&\mid \; \text{<>} \\
48.&\quad &&\mid \; \text{=} \\
49.&\quad addop &&::= \text{+} \\
50.&\quad &&\mid \; \text{-} \\
51.&\quad &&\mid \; \text{or} \\
52.&\quad mulop &&::= \text{*} \\
53.&\quad &&\mid \; \text{/} \\
54.&\quad &&\mid \; \text{\%} \\
55.&\quad &&\mid \; \text{and}
\intertext{\subsection{Regras de Formação de Tokens}}
56.&\quad constant &&::= \mathit{integer\_const} \; \mathit{parte\_real} \\
57.&\quad &&\mid \; literal \\
58.&\quad \mathit{parte\_real} &&::= \text{.} \; digit^{+} \\
59.&\quad &&\mid \; \lambda \\
60.&\quad \mathit{integer\_const} &&\to nonzero \; digit^{*} \\
61.&\quad &&\mid \; \text{0} \\
62.&\quad literal &&\to \text{"} \; caractere^{*} \; \text{"} \\
63.&\quad identifier &&\to letter \; (letter \mid digit)^{*} \\
64.&\quad letter &&\to \text{[A-Za-z]} \\
65.&\quad digit &&\to \text{[0-9]} \\
66.&\quad nonzero &&\to \text{[1-9]} \\
67.&\quad caractere &&\to \text{um dos 256 caracteres do conjunto} \\
&\quad &&\quad \text{ASCII, exceto aspas e quebra de linha}
\end{align*}

\subsection{Modificações na Gramática para Adequação LL(1)}

A gramática original fornecida foi modificada utilizando duas técnicas: eliminação de
recursão à esquerda e fatoração à esquerda.

\subsubsection{Eliminação de Recursão à Esquerda}

\paragraph{Expressões:}
\begin{align*}
expression      &\to simple\!-\!expr \; expression' \\
expression'     &\to relop \; simple\!-\!expr \mid \lambda
\end{align*}

\paragraph{Expressões Simples:}
\begin{align*}
simple\!-\!expr    &\to term \; simple\!-\!expr' \\
simple\!-\!expr'   &\to addop \; term \; simple\!-\!expr' \mid \lambda
\end{align*}

Na implementação, \texttt{simpleExprPrime()} utiliza iteração (\texttt{while}) em vez de recursão, por
ser equivalente para recursão.

\paragraph{Termos:}
\begin{align*}
term    &\to factor\!-\!a \; term' \\
term'   &\to mulop \; factor\!-\!a \; term' \mid \lambda
\end{align*}

Analogamente, \texttt{termPrime()} utiliza iteração.

\subsubsection{Fatoração à Esquerda}

O comando \texttt{if} possuía duas produções com prefixo comum:
\begin{align*}
if\!-\!stmt     &\to \text{if } \text{(} \; condition \; \text{)} \; \text{\{} \; stmt\!-\!list \; \text{\}} \; if\!-\!stmt' \\
if\!-\!stmt'    &\to \text{else } \text{\{} \; stmt\!-\!list \; \text{\}} \mid \lambda
\end{align*}

O método correspondente é \texttt{ifStmtPrime()} no código.

A constante também possuía duas produções com prefixo comum:
\begin{align*}
constant        &\to integer\_const \mid literal \mid real\_const \\
real\_const     &\to integer\_const \; \text{"."} \; digit^{+}
\end{align*}
As alternativas \texttt{integer\_const} e \texttt{real\_const} compartilham o prefixo \texttt{integer\_const}, impossibilitando a decisão determinística para o parser $LL(1)$. Aplicou-se a fatoração à esquerda, introduzindo \texttt{parte\_real}:
\begin{align*}
constant        &\to integer\_const \; parte\_real \mid literal \\
parte\_real     &\to \text{"."} \; digit^{+} \mid \lambda
\end{align*}
Agora, após reconhecer \texttt{integer\_const}, o parser consulta o lookahead: se for \texttt{"."}, escolhe \texttt{parte\_real \to "." digit^{+}} (constante real); caso contrário, escolhe $\lambda$ (constante inteira).


\subsubsection{Observações sobre a implementação}

\normalfont\normalsize

Na implementação real, as produções com $*$, $+$ e listas ($decl\!-\!list$, $ident\!-\!list'$,
$stmt\!-\!list$, $simple\!-\!expr'$, $term'$) foram implementadas com laços \texttt{while} em vez de recursão, por
questão de eficiência e para evitar estouro de pilha em entradas grandes. A semântica é
idêntica à das gramáticas recursivas apresentadas.

O método \texttt{constant()} foi separado de \texttt{factor()}:
\begin{align*}
factor          &\to identifier \mid constant \mid \text{(} \; expression \; \text{)}
\end{align*}


E \texttt{factorA()} adiciona os prefixos opcionais:
\begin{align*}
factor\!-\!a    &\to \text{not } factor \mid \text{-} factor \mid factor
\end{align*}

\section{Conjuntos FIRST e FOLLOW}

\begingroup
\centering
\begin{longtable}{|l|p{6cm}|p{6cm}|}
\hline
\textbf{Não-terminal} & \textbf{FIRST} & \textbf{FOLLOW} \\ \hline
\endfirsthead
\hline
\textbf{Não-terminal} & \textbf{FIRST} & \textbf{FOLLOW} \\ \hline
\endhead
\hline
\endfoot
\endlastfoot

program         & class                                  & \$ \\ \hline
decl-list       & int string float $\lambda$              & \{ \\ \hline
decl            & int string float                        & ; \\ \hline
ident-list      & $[A-Za-z]$                              & ; \\ \hline
type            & int string float                        & $[A-Za-z]$ \\ \hline
body            & \{                                      & \} \\ \hline
stmt-list       & $[A-Za-z]$ if do repeat read write      & \} \\ \hline
stmt            & $[A-Za-z]$ if do repeat read write      & ; \\ \hline
assign-stmt     & $[A-Za-z]$                              & ; \\ \hline
if-stmt         & if                                       & ; \\ \hline
if-stmt'        & else $\lambda$                           & ; \\ \hline
do-stmt         & do                                       & ; \\ \hline
do-suffix       & while                                    & ; \\ \hline
repeat-stmt     & repeat                                   & ; \\ \hline
stmt-suffix     & until                                    & ; \\ \hline
read-stmt       & read                                     & ; \\ \hline
write-stmt      & write                                    & ; \\ \hline
writable        & $[A-Za-z]$ $[1-9]$ 0 " ( not -          & ) \\ \hline
condition       & $[A-Za-z]$ $[1-9]$ 0 " ( not -          & ) \\ \hline
expression      & $[A-Za-z]$ $[1-9]$ 0 " ( not -          & ) \\ \hline
expression'     & > >= < <= <> = $\lambda$                 & ) \\ \hline
simple-expr     & $[A-Za-z]$ $[1-9]$ 0 " ( not -          & ) > >= < <= <> = ; \\ \hline
simple-expr'    & or + - $\lambda$                         & ) > >= < <= <> = ; \\ \hline
term            & $[A-Za-z]$ $[1-9]$ 0 " ( not -          & or + - ) > >= < <= <> = ; \\ \hline
term'           & * / \% and $\lambda$                     & or + - ) > >= < <= <> = ; \\ \hline
factor-a        & $[A-Za-z]$ $[1-9]$ 0 " ( not -          & * / \% and or + - ) > >= < <= <> = ; \\ \hline
factor          & $[A-Za-z]$ $[1-9]$ 0 " (                & * / \% and or + - ) > >= < <= <> = ; \\ \hline
relop           & > >= < <= <> =                           & $[A-Za-z]$ $[1-9]$ 0 " ( not - \\ \hline
addop           & or + -                                   & $[A-Za-z]$ $[1-9]$ 0 " ( not - \\ \hline
mulop           & * / \% and                               & $[A-Za-z]$ $[1-9]$ 0 " ( not - \\ \hline
constant        & $[1-9]$ 0 "                              & * / \% and or + - ) > >= < <= <> = ; \\ \hline
parte\_real     & . $\lambda$                              & * / \% and or + - ) > >= < <= <> = ; \\ \hline
integer\_const  & $[1-9]$ 0                                & . * / \% and or + - ) > >= < <= <> = ; \\ \hline
literal         & "                                        & * / \% and or + - ) > >= < <= <> = ; \\ \hline
identifier      & $[A-Za-z]$                              & \{ ; , \} := ) * / \% and or + - > >= < <= <> = \\ \hline
letter          & $[A-Za-z]$                              & \{ ; , \} := ) * / \% and or + - > >= < <= <> = $[A-Za-z]$ $[0-9]$ \\ \hline
digit           & $[0-9]$                                 & * / \% and or + - ) > >= < <= <> = $[0-9]$ . \{ ; , \} := $[A-Za-z]$ \\ \hline
nonzero         & $[1-9]$                                 & . * / \% and or + - ) > >= < <= <> = $[0-9]$ \\ \hline
caractere       & ASCII, exceto aspas e /n                & " \\ \hline
\end{longtable}
\endgroup

\normalfont\normalsize

\section{Estrutura da Implementação Sintática}

O parser é estruturado em torno de três funções auxiliares:

\begin{itemize}
    \item \texttt{advance()} --- chama \texttt{lex.scan()} para avançar o lookahead, tratando exceções de leitura.
    \item \texttt{match(int tag)} --- verifica se o token corrente corresponde ao esperado; em caso positivo, avança; em caso negativo, chama \texttt{error()}.
    \item \texttt{error(String msg)} --- exibe a linha (\texttt{Lexer.line}) e a mensagem, e encerra com \texttt{System.exit(1)}.
\end{itemize}

O ponto de entrada é o método público \texttt{program()}, chamado diretamente pela \texttt{Main}:

\begin{lstlisting}[language=Java, caption=Classe principal (Main.java)]
import lexico.Lexer;
import sintatico.Parser;

public class Main {
    public static void main(String[] args) {
        try {
            Lexer lex = new Lexer(args[0]);
            Parser parser = new Parser(lex);
            parser.program();
            System.out.println("Programa sintaticamente correto.");
        } catch (FileNotFoundException e) {
            System.err.println("Arquivo nao encontrado.");
        } catch (IOException e) {
            System.err.println("Erro de leitura: " + e.getMessage());
        }
    }
}
\end{lstlisting}

\begin{lstlisting}[language=Java, caption=Esqueleto do Parser]
package sintatico;

import lexico.*;

public class Parser {
    private Lexer lex;
    private Token lookahead;

    public Parser(Lexer lex) {
        this.lex = lex;
        advance();
    }

    private void advance() {
        try {
            lookahead = lex.scan();
        } catch (Exception e) {
            error("Erro de leitura: " + e.getMessage());
        }
    }

    private void match(int tag) {
        if (lookahead.tag == tag) {
            advance();
        } else {
            error("Esperado " + tagToString(tag)
                + ", encontrado " + tokenToString(lookahead));
        }
    }

    private void error(String msg) {
        System.err.println("Erro sintatico na linha "
            + Lexer.line + ": " + msg);
        System.exit(1);
    }

    public void program() {
        match(Tag.CLASS);
        match(Tag.ID);
        match('{');
        if (isDeclListStart()) declList();
        body();
        match('}');
        if (lookahead.tag != -1)
            error("Conteudo extra apos o fim do programa");
    }

    // if_stmt ::= if ( condition ) { stmt_list } if_stmt'
    private void ifStmt() {
        match(Tag.IF);
        match('(');
        condition();
        match(')');
        match('{');
        stmtList();
        match('}');
        ifStmtPrime();
    }

    // if_stmt' ::= else { stmt_list } | epsilon
    private void ifStmtPrime() {
        if (lookahead.tag == Tag.ELSE) {
            match(Tag.ELSE);
            match('{');
            stmtList();
            match('}');
        }
    }
}
\end{lstlisting}

\normalfont\normalsize

\section{Casos de Teste (7 a 12)}

Conforme as diretrizes da Etapa 2, os testes 7 a 12 são versões dos testes 1 a 6 com ajustes incrementais para eliminar erros léxicos e sintáticos. Para cada teste, executou-se o parser, identificou-se o primeiro erro, corrigiu-se a entrada e reexecutou-se até a aprovação completa.

\subsection{Teste 7 --- Cálculo de área}

\textbf{1º erro -- identificador começando com dígito}

\begin{lstlisting}[language=Java]
class Teste1{
    int base, 1altura;
    float area;

    write("Digite o valor da base:");
    read (base);
    write("Digite o valor da altura:");
    read (altura;
    area :=base*altura/2.0
    write("A area e: " + area);
}
\end{lstlisting}

Código 5: Teste7.txt -- com erros

\begin{lstlisting}[language=bash]
> java -cp out Main compilador\Teste7.txt
Erro sintatico na linha 2: Esperado identificador, encontrado numero 1
\end{lstlisting}

Análise: O identificador \texttt{1altura} inicia com dígito. O lexer tokeniza \texttt{1} como número e \texttt{altura} como identificador. Após a vírgula em \texttt{base,}, o parser espera um identificador para a lista de declarações, mas encontra o número \texttt{1}.

Correção: \texttt{1altura} $\to$ \texttt{altura1}

\textbf{2º erro -- corpo sem chaves de abertura e fechamento}

\begin{lstlisting}[language=Java]
class Teste1{
    int base, altura1;
    float area;

    write("Digite o valor da base:");
    read (base);
    write("Digite o valor da altura:");
    read (altura;
    area :=base*altura/2.0
    write("A area e: " + area);
}
\end{lstlisting}

Código 5.1: Teste7\_corrigido1.txt

\begin{lstlisting}[language=bash]
> java -cp out Main compilador\Teste7_corrigido1.txt
Erro sintatico na linha 5: Esperado '{', encontrado 'write'
\end{lstlisting}

Análise: Após as declarações de variáveis, o parser espera \texttt{\{} para iniciar o corpo (\texttt{body}) da classe, mas encontra \texttt{write}.

Correção: Adicionado par de chaves \texttt{\{} \texttt{\}} (abertura e fechamento) para o corpo

\textbf{3º erro -- parêntese não fechado}

\begin{lstlisting}[language=Java]
class Teste1{
    int base, altura1;
    float area;
    {
    write("Digite o valor da base:");
    read (base);
    write("Digite o valor da altura:");
    read (altura;
    area :=base*altura/2.0
    write("A area e: " + area);
    }
}
\end{lstlisting}

Código 5.2: Teste7\_corrigido2.txt

\begin{lstlisting}[language=bash]
> java -cp out Main compilador\Teste7_corrigido2.txt
Erro sintatico na linha 8: Esperado ')', encontrado ';'
\end{lstlisting}

Análise: O comando \texttt{read (altura;} não possui o parêntese de fechamento. O parser encontra \texttt{;} e esperava \texttt{)}.

Correção: \texttt{read(altura;} $\to$ \texttt{read(altura);}

\textbf{4º erro -- ponto e vírgula ausente}

\begin{lstlisting}[language=Java]
class Teste1{
    int base, altura1;
    float area;
    {
    write("Digite o valor da base:");
    read (base);
    write("Digite o valor da altura:");
    read (altura);
    area :=base*altura/2.0
    write("A area e: " + area);
    }
}
\end{lstlisting}

Código 5.3: Teste7\_corrigido3.txt

\begin{lstlisting}[language=bash]
> java -cp out Main compilador\Teste7_corrigido3.txt
Erro sintatico na linha 10: Esperado ';', encontrado 'write'
\end{lstlisting}

Análise: A atribuição \texttt{area := base * altura / 2.0} não possui \texttt{;} ao final. O parser encontra \texttt{write} na linha seguinte.

Correção: Adicionado \texttt{;} ao final de \texttt{area := base * altura / 2.0}

\textbf{Teste aprovado}

\begin{lstlisting}[language=Java]
class Teste1{
    int base, altura1;
    float area;
    {
    write("Digite o valor da base:");
    read (base);
    write("Digite o valor da altura:");
    read (altura);
    area :=base*altura/2.0;
    write("A area e: " + area);
    }
}
\end{lstlisting}

Código 6: Teste7\_corrigido.txt

\begin{lstlisting}[language=bash]
> java -cp out Main compilador\Teste7_corrigido.txt
Programa sintaticamente correto.
\end{lstlisting}

\subsection{Teste 8 --- Comentários e if-else}

\textbf{1º erro -- corpo sem chaves de abertura e fechamento}

\begin{lstlisting}[language=Java]
class Teste2{
    /* Teste de comentario
    com mais de uma linha */
    float a;

    write("Entre com o valor de a: ");
    read (a);

    if (a%2 = 0){
        write (a+ "e par").
    }
    else{
        write (a+ "e impar");
    }
}
\end{lstlisting}

Código 7: Teste8.txt -- com erros

\begin{lstlisting}[language=bash]
> java -cp out Main compilador\Teste8.txt
Erro sintatico na linha 6: Esperado '{', encontrado 'write'
\end{lstlisting}

Análise: Após a declaração \texttt{float a;}, o corpo (\texttt{body}) espera \texttt{\{} para iniciar a lista de comandos, mas encontra \texttt{write}.

Correção: Adicionado par de chaves \texttt{\{} \texttt{\}} (abertura e fechamento) para o corpo, antes e depois dos comandos

\textbf{2º erro -- ponto final no lugar de ponto e vírgula}

\begin{lstlisting}[language=Java]
class Teste2{
    /* Teste de comentario
    com mais de uma linha */
    float a;
    {
    write("Entre com o valor de a: ");
    read (a);

    if (a%2 = 0){
        write (a+ "e par").
    }
    else{
        write (a+ "e impar");
    }
    }
}
\end{lstlisting}

Código 7.1: Teste8\_corrigido1.txt

\begin{lstlisting}[language=bash]
> java -cp out Main compilador\Teste8_corrigido1.txt
Erro sintatico na linha 10: Esperado ';', encontrado '.'
\end{lstlisting}

Análise: O comando \texttt{write (a+ "e par")} termina com ponto final (\texttt{.}) em vez de ponto e vírgula (\texttt{;}).

Correção: \texttt{write (a+ "e par").} $\to$ \texttt{write (a+ "e par");}

\textbf{3º erro -- ponto e vírgula ausente após if-else}

\begin{lstlisting}[language=Java]
class Teste2{
    /* Teste de comentario
    com mais de uma linha */
    float a;
    {
    write("Entre com o valor de a: ");
    read (a);

    if (a%2 = 0){
        write (a+ "e par");
    }
    else{
        write (a+ "e impar");
    }
    }
}
\end{lstlisting}

Código 7.2: Teste8\_corrigido2.txt

\begin{lstlisting}[language=bash]
> java -cp out Main compilador\Teste8_corrigido2.txt
Erro sintatico na linha 15: Esperado ';', encontrado '}'
\end{lstlisting}

Análise: A gramática exige \texttt{;} após todo comando na \texttt{stmt-list}, incluindo o comando \texttt{if-else}. O parser encontra \texttt{\}} de fechamento do corpo e esperava \texttt{;} para finalizar o if-else.

Correção: Adicionado \texttt{;} após o \texttt{\}} de fechamento do else

\textbf{Teste aprovado}

\begin{lstlisting}[language=Java]
class Teste2{
    /* Teste de comentario
    com mais de uma linha */
    float a;
    {
    write("Entre com o valor de a: ");
    read (a);

    if (a%2 = 0){
        write (a+ "e par");
    }
    else{
        write (a+ "e impar");
    };
    }
}
\end{lstlisting}

Código 8: Teste8\_corrigido.txt

\begin{lstlisting}[language=bash]
> java -cp out Main compilador\Teste8_corrigido.txt
Programa sintaticamente correto.
\end{lstlisting}

\subsection{Teste 9 --- If aninhado}

\textbf{1º erro -- corpo sem chaves de abertura e fechamento}

\begin{lstlisting}[language=Java]
class MinhaClasse
    {
    float a, b, c;

        write("Digite um numero");
        read(a);
        write("Digite outro numero: ");
        read(b);
        write("Digite mais um numero: );
        read(c;

        maior := 0;

        if ( a>b && a>c )
            maior = a;
        else
            if (b>c)
                maior = b;
            else
                maior = c;

        write("O maior numero e: ");
        write(maior);
\end{lstlisting}

Código 7: Teste9.txt -- com erros

\begin{lstlisting}[language=bash]
> java -cp out Main compilador\Teste9.txt
Erro sintatico na linha 5: Esperado '{', encontrado 'write'
\end{lstlisting}

Análise: O \texttt{\{} na linha 2 é a abertura da classe (após \texttt{class MinhaClasse}). Após a declaração \texttt{float a, b, c;}, o corpo espera \texttt{\{} mas encontra \texttt{write}.

Correção: Adicionado par de chaves \texttt{\{} \texttt{\}} (abertura e fechamento) para o corpo e fechado a string literal na linha 9

\textbf{2º erro -- parêntese não fechado}

\begin{lstlisting}[language=Java]
class MinhaClasse
    {
    float a, b, c;
    {
        write("Digite um numero");
        read(a);
        write("Digite outro numero: ");
        read(b);
        write("Digite mais um numero: ");
        read(c;

        maior := 0;

        if ( a>b && a>c )
            maior = a;
        else
            if (b>c)
                maior = b;
            else
                maior = c;

        write("O maior numero e: ");
        write(maior);
    }
}
\end{lstlisting}

Código 7.1: Teste9\_corrigido1.txt

\begin{lstlisting}[language=bash]
> java -cp out Main compilador\Teste9_corrigido1.txt
Erro sintatico na linha 10: Esperado ')', encontrado ';'
\end{lstlisting}

Análise: \texttt{read(c;} não possui o parêntese de fechamento \texttt{)}.

Correção: \texttt{read(c;} $\to$ \texttt{read(c);}

\textbf{3º erro -- operador \texttt{\&\&} (primeira parte: substituir \texttt{\&\&} por \texttt{and})}

\begin{lstlisting}[language=Java]
class MinhaClasse
    {
    float a, b, c;
    {
        write("Digite um numero");
        read(a);
        write("Digite outro numero: ");
        read(b);
        write("Digite mais um numero: ");
        read(c);

        maior := 0;

        if ( a>b && a>c )
            maior = a;
        else
            if (b>c)
                maior = b;
            else
                maior = c;

        write("O maior numero e: ");
        write(maior);
    }
}
\end{lstlisting}

Código 7.2: Teste9\_corrigido2.txt

\begin{lstlisting}[language=bash]
> java -cp out Main compilador\Teste9_corrigido2.txt
Erro sintatico na linha 14: Esperado ')', encontrado '&'
\end{lstlisting}

Análise: O operador \texttt{\&\&} não pertence à gramática.

Correção: Substituir \texttt{\&\&} por \texttt{and}

\textbf{4º erro -- operador \texttt{and} (segunda parte: adicionar parênteses nas expressões)}

\begin{lstlisting}[language=Java]
class MinhaClasse
    {
    float a, b, c;
    {
        write("Digite um numero");
        read(a);
        write("Digite outro numero: ");
        read(b);
        write("Digite mais um numero: ");
        read(c);

        maior := 0;

        if ( a>b and a>c )
            maior = a;
        else
            if (b>c)
                maior = b;
            else
                maior = c;

        write("O maior numero e: ");
        write(maior);
    }
}
\end{lstlisting}

Código 7.3: Teste9\_corrigido3a.txt

\begin{lstlisting}[language=bash]
> java -cp out Main compilador\Teste9_corrigido3a.txt
Erro sintatico na linha 14: Esperado ')', encontrado '>'
\end{lstlisting}

Análise: A simples substitui\c{c}\~ao de \texttt{\&\&} por \texttt{and} n\~ao \'e suficiente, pois \texttt{and} \'e um operador multiplicativo (mulop) na gram\'atica. A express\~ao \texttt{a>b and a>c} \'e interpretada como \texttt{a > (b and a) > c}. O analisador consome \texttt{a > b and a} e, ao encontrar o segundo \texttt{>}, espera o fechamento do par\^entese da condi\c{c}\~ao do if.

Correção: Adicionar par\^enteses em cada express\~ao relacional: \texttt{(a>b)} e \texttt{(a>c)}

\textbf{5º erro -- if externo sem chaves no corpo}

\begin{lstlisting}[language=Java]
class MinhaClasse
    {
    float a, b, c;
    {
        write("Digite um numero");
        read(a);
        write("Digite outro numero: ");
        read(b);
        write("Digite mais um numero: ");
        read(c);

        maior := 0;

        if ((a > b) and (a > c))
            maior = a;
        else
            if (b>c)
                maior = b;
            else
                maior = c;

        write("O maior numero e: ");
        write(maior);
    }
}
\end{lstlisting}

Código 7.4: Teste9\_corrigido4b.txt

\begin{lstlisting}[language=bash]
> java -cp out Main compilador\Teste9_corrigido4b.txt
Erro sintatico na linha 15: Esperado '{', encontrado 'maior'
\end{lstlisting}


Análise: O if externo n\~ao possui chaves para delimitar seu corpo. A gram\'atica exige \texttt{\{} ap\'os a condi\c{c}\~ao do \texttt{if} para iniciar o bloco de comandos.

Corre\c{c}\~ao: Adicionar \texttt{\{} e \texttt{\}} ao redor de \texttt{maior = a;} (corpo do if externo)

\textbf{6º erro -- atribuição com \texttt{=} no if externo}

\begin{lstlisting}[language=Java]
class MinhaClasse
    {
    float a, b, c;
    {
        write("Digite um numero");
        read(a);
        write("Digite outro numero: ");
        read(b);
        write("Digite mais um numero: ");
        read(c);

        maior := 0;

        if ((a > b) and (a > c)){
            maior = a;
        }
        else
            if (b>c)
                maior = b;
            else
                maior = c;

        write("O maior numero e: ");
        write(maior);
    }
}
\end{lstlisting}

Código 7.5: Teste9\_corrigido5.txt

\begin{lstlisting}[language=bash]
> java -cp out Main compilador\Teste9_corrigido5.txt
Erro sintatico na linha 15: Esperado :=, encontrado '='
\end{lstlisting}


Análise: O if externo agora possui chaves, mas a atribuição \texttt{maior = a;} utiliza \texttt{=} em vez de \texttt{:=}.

Correção: \texttt{maior = a;} $\to$ \texttt{maior := a;}

\textbf{7º erro -- else externo sem chaves}

\begin{lstlisting}[language=Java]
class MinhaClasse
    {
    float a, b, c;
    {
        write("Digite um numero");
        read(a);
        write("Digite outro numero: ");
        read(b);
        write("Digite mais um numero: ");
        read(c);

        maior := 0;

        if ((a > b) and (a > c)){
            maior := a;
        }
        else
            if (b>c)
                maior = b;
            else
                maior = c;

        write("O maior numero e: ");
        write(maior);
    }
}
\end{lstlisting}

Código 7.6: Teste9\_corrigido6.txt

\begin{lstlisting}[language=bash]
> java -cp out Main compilador\Teste9_corrigido6.txt
Erro sintatico na linha 18: Esperado '{', encontrado 'if'
\end{lstlisting}


Análise: O else externo n\~ao possui chaves. A gram\'atica exige \texttt{else \{ stmt-list \}}.

Corre\c{c}\~ao: Adicionar \texttt{\{} e \texttt{\}} ao redor do corpo do else externo

\textbf{8º erro -- if interno sem chaves no corpo}

\begin{lstlisting}[language=Java]
class MinhaClasse
    {
    float a, b, c;
    {
        write("Digite um numero");
        read(a);
        write("Digite outro numero: ");
        read(b);
        write("Digite mais um numero: ");
        read(c);

        maior := 0;

        if ((a > b) and (a > c)){
            maior := a;
        }
        else{
            if (b>c)
                maior = b;
            else
                maior = c;
        }

        write("O maior numero e: ");
        write(maior);
    }
}
\end{lstlisting}

Código 7.7: Teste9\_corrigido7.txt

\begin{lstlisting}[language=bash]
> java -cp out Main compilador\Teste9_corrigido7.txt
Erro sintatico na linha 19: Esperado '{', encontrado 'maior'
\end{lstlisting}


Análise: O if interno (\texttt{if (b>c)}) também precisa de chaves para seu corpo.

Correção: Adicionar \texttt{\{} e \texttt{\}} ao redor de \texttt{maior = b;}

\textbf{9º erro -- atribuição com \texttt{=} no if interno}

\begin{lstlisting}[language=Java]
class MinhaClasse
    {
    float a, b, c;
    {
        write("Digite um numero");
        read(a);
        write("Digite outro numero: ");
        read(b);
        write("Digite mais um numero: ");
        read(c);

        maior := 0;

        if ((a > b) and (a > c)){
            maior := a;
        }
        else{
            if (b>c){
                maior = b;
            }
            else
                maior = c;
        }

        write("O maior numero e: ");
        write(maior);
    }
}
\end{lstlisting}

Código 7.8: Teste9\_corrigido8.txt

\begin{lstlisting}[language=bash]
> java -cp out Main compilador\Teste9_corrigido8.txt
Erro sintatico na linha 19: Esperado :=, encontrado '='
\end{lstlisting}


Análise: O if interno agora possui chaves, mas a atribuição \texttt{maior = b;} utiliza \texttt{=} em vez de \texttt{:=}.

Correção: \texttt{maior = b;} $\to$ \texttt{maior := b;}

\textbf{10º erro -- else interno sem chaves}

\begin{lstlisting}[language=Java]
class MinhaClasse
    {
    float a, b, c;
    {
        write("Digite um numero");
        read(a);
        write("Digite outro numero: ");
        read(b);
        write("Digite mais um numero: ");
        read(c);

        maior := 0;

        if ((a > b) and (a > c)){
            maior := a;
        }
        else{
            if (b>c){
                maior := b;
            }
            else
                maior = c;
        }

        write("O maior numero e: ");
        write(maior);
    }
}
\end{lstlisting}

Código 7.9: Teste9\_corrigido9.txt

\begin{lstlisting}[language=bash]
> java -cp out Main compilador\Teste9_corrigido9.txt
Erro sintatico na linha 22: Esperado '{', encontrado 'maior'
\end{lstlisting}
Análise: O else interno n\~ao possui chaves para delimitar seu corpo.

Corre\c{c}\~ao: Adicionar \texttt{\{} e \texttt{\}} ao redor de \texttt{maior = c;}

\textbf{11º erro -- atribuição com \texttt{=} no else interno}

\begin{lstlisting}[language=Java]
class MinhaClasse
    {
    float a, b, c;
    {
        write("Digite um numero");
        read(a);
        write("Digite outro numero: ");
        read(b);
        write("Digite mais um numero: ");
        read(c);

        maior := 0;

        if ((a > b) and (a > c)){
            maior := a;
        }
        else{
            if (b>c){
                maior := b;
            }
            else{
                maior = c;
            }
        }

        write("O maior numero e: ");
        write(maior);
    }
}
\end{lstlisting}

Código 7.10: Teste9\_corrigido10.txt

\begin{lstlisting}[language=bash]
> java -cp out Main compilador\Teste9_corrigido10.txt
Erro sintatico na linha 22: Esperado :=, encontrado '='
\end{lstlisting}

Análise: O else interno agora possui chaves, mas a atribuição \texttt{maior = c;} utiliza \texttt{=} em vez de \texttt{:=}.

Correção: \texttt{maior = c;} $\to$ \texttt{maior := c;}

\textbf{12º erro -- ponto e vírgula ausente no if-else interno}

\begin{lstlisting}[language=Java]
class MinhaClasse
    {
    float a, b, c;
    {
        write("Digite um numero");
        read(a);
        write("Digite outro numero: ");
        read(b);
        write("Digite mais um numero: ");
        read(c);

        maior := 0;

        if ((a > b) and (a > c)){
            maior := a;
        }
        else{
            if (b>c){
                maior := b;
            }
            else{
                maior := c;
            }
        }

        write("O maior numero e: ");
        write(maior);
    }
}
\end{lstlisting}

Código 7.11: Teste9\_corrigido11.txt

\begin{lstlisting}[language=bash]
> java -cp out Main compilador\Teste9_corrigido11.txt
Erro sintatico na linha 24: Esperado ';', encontrado '}'
\end{lstlisting}

Análise: O if-else interno é um comando na \texttt{stmt-list} e deve terminar com \texttt{;}.

Correção: Adicionado \texttt{;} após o \texttt{\}} de fechamento do else interno

\textbf{13º erro -- ponto e vírgula ausente no if-else externo}

\begin{lstlisting}[language=Java]
class MinhaClasse
    {
    float a, b, c;
    {
        write("Digite um numero");
        read(a);
        write("Digite outro numero: ");
        read(b);
        write("Digite mais um numero: ");
        read(c);

        maior := 0;

        if ((a > b) and (a > c)){
            maior := a;
        }
        else{
            if (b>c){
                maior := b;
            }
            else{
                maior := c;
            };
        }

        write("O maior numero e: ");
        write(maior);
    }
}
\end{lstlisting}

Código 7.12: Teste9\_corrigido12.txt

\begin{lstlisting}[language=bash]
> java -cp out Main compilador\Teste9_corrigido12.txt
Erro sintatico na linha 26: Esperado ';', encontrado 'write'
\end{lstlisting}

Análise: O if-else externo também precisa de \texttt{;} para terminar o comando.

Correção: Adicionado \texttt{;} após o \texttt{\}} de fechamento do else externo

\textbf{Teste aprovado}

\begin{lstlisting}[language=Java]
class MinhaClasse
    {
    float a, b, c;
    {
        write("Digite um numero");
        read(a);
        write("Digite outro numero: ");
        read(b);
        write("Digite mais um numero: ");
        read(c);

        maior := 0;

        if ((a > b) and (a > c)){
            maior := a;
        }
        else{
            if (b > c){
                maior := b;
            }
            else{
                maior := c;
            };
        };
        write("O maior numero e: ");
        write(maior);
    }
}
\end{lstlisting}

Código 8: Teste9\_corrigido.txt

\begin{lstlisting}[language=bash]
> java -cp out Main compilador\Teste9_corrigido.txt
Programa sintaticamente correto.
\end{lstlisting}

\subsection{Teste 10 --- Repeat-until}

\textbf{1º erro -- programa sem a palavra-chave \texttt{class}}

\begin{lstlisting}[language=Java]
{
    // Outro programa de teste
    int idade;
    string nome, texto;
    float bonus, salario;

    repeat{
        write("Digite o seu nome: ");
        read(nome);
        write("Digite o seu sobrenome");
        read(sobrenome);
        write("Digite a sua idade: ");
        read (idade);
        write("Digite o salario: ");
        read (salrio);
        write("Digite o bonus: ");
        read (bonus);
        if (idade>0 && salario > 0 and bonus >= 0){
            salarioLiquido = +salario salario*0.11) + (bonus
            109.8);
            texto = nome + " " + sobrenome + ":";
            write (text);
            write("Salario liquido: ");
            write(salario);
        };
    {until (idade <=0 or Salario <=0 or bonus <0);
}
\end{lstlisting}

Código 9: Teste10.txt -- com erros

\begin{lstlisting}[language=bash]
> java -cp out Main compilador\Teste10.txt
Erro sintatico na linha 1: Esperado class, encontrado '{'
\end{lstlisting}

Análise: O programa inicia diretamente com \texttt{\{} sem a palavra reservada \texttt{class}.

Correção: Adicionado \texttt{class Teste4\{} antes do \texttt{\{} de abertura

\textbf{2º erro -- corpo do programa sem chaves de abertura e fechamento}

\begin{lstlisting}[language=Java]
class Teste4{
    // Outro programa de teste
    int idade;
    string nome, texto;
    float bonus, salario;

    repeat{
        write("Digite o seu nome: ");
        read(nome);
        write("Digite o seu sobrenome");
        read(sobrenome);
        write("Digite a sua idade: ");
        read (idade);
        write("Digite o salario: ");
        read (salrio);
        write("Digite o bonus: ");
        read (bonus);
        if (idade>0 && salario > 0 and bonus >= 0){
            salarioLiquido = +salario salario*0.11) + (bonus
            109.8);
            texto = nome + " " + sobrenome + ":";
            write (text);
            write("Salario liquido: ");
            write(salario);
        };
    {until (idade <=0 or Salario <=0 or bonus <0);
    }
}
}
\end{lstlisting}

Código 9.1: Teste10\_corrigido1.txt

\begin{lstlisting}[language=bash]
> java -cp out Main compilador\Teste10_corrigido1.txt
Erro sintatico na linha 7: Esperado '{', encontrado 'repeat'
\end{lstlisting}

Análise: Após as declarações (\texttt{int}, \texttt{string}, \texttt{float}), o parser espera \texttt{\{} para iniciar o corpo (\texttt{body}), mas encontra \texttt{repeat}.

Correção: Adicionado par de chaves \texttt{\{} \texttt{\}} (abertura e fechamento) para o corpo, antes e depois dos comandos

\textbf{3º erro -- operador \texttt{\&\&} (primeira parte: substituir \texttt{\&\&} por \texttt{and})}

\begin{lstlisting}[language=Java]
class Teste4{
    // Outro programa de teste
    int idade;
    string nome, texto;
    float bonus, salario;
    {
    repeat{
        write("Digite o seu nome: ");
        read(nome);
        write("Digite o seu sobrenome");
        read(sobrenome);
        write("Digite a sua idade: ");
        read (idade);
        write("Digite o salario: ");
        read (salrio);
        write("Digite o bonus: ");
        read (bonus);
        if (idade>0 && salario > 0 and bonus >= 0){
            salarioLiquido = +salario salario*0.11) + (bonus
            109.8);
            texto = nome + " " + sobrenome + ":";
            write (text);
            write("Salario liquido: ");
            write(salario);
        };
    {until (idade <=0 or Salario <=0 or bonus <0);
    }
}
}
\end{lstlisting}

Código 9.2: Teste10\_corrigido2.txt

\begin{lstlisting}[language=bash]
> java -cp out Main compilador\Teste10_corrigido2.txt
Erro sintatico na linha 18: Esperado ')', encontrado '&'
\end{lstlisting}

Análise: O operador \texttt{\&\&} não pertence à gramática.

Correção: Substituir \texttt{\&\&} por \texttt{and}

\textbf{4º erro -- operador \texttt{and} (segunda parte: adicionar parênteses nas expressões)}

\begin{lstlisting}[language=Java]
class Teste4{
    // Outro programa de teste
    int idade;
    string nome, texto;
    float bonus, salario;
    {
    repeat{
        write("Digite o seu nome: ");
        read(nome);
        write("Digite o seu sobrenome");
        read(sobrenome);
        write("Digite a sua idade: ");
        read (idade);
        write("Digite o salario: ");
        read (salrio);
        write("Digite o bonus: ");
        read (bonus);
        if (idade>0 and salario > 0 and bonus >= 0){
            salarioLiquido = +salario salario*0.11) + (bonus
            109.8);
            texto = nome + " " + sobrenome + ":";
            write (text);
            write("Salario liquido: ");
            write(salario);
        };
    {until (idade <=0 or Salario <=0 or bonus <0);
    }
}
}
\end{lstlisting}

Código 9.3: Teste10\_corrigido3a.txt

\begin{lstlisting}[language=bash]
> java -cp out Main compilador\Teste10_corrigido3a.txt
Erro sintatico na linha 18: Esperado ')', encontrado '>'
\end{lstlisting}

Análise: A simples substituição de \texttt{\&\&} por \texttt{and} não é suficiente, pois \texttt{and} é um operador multiplicativo (mulop) na gramática. A expressão \texttt{idade>0 and salario > 0} é interpretada como \texttt{idade > (0 and salario) > 0}, portanto o analisador espera \texttt{)} para fechar a condição do if.

Correção: Adicionar parênteses em cada expressão relacional: \texttt{idade>0} $\to$ \texttt{(idade>0)}, \texttt{salario > 0} $\to$ \texttt{(salario > 0)}, \texttt{bonus >= 0} $\to$ \texttt{(bonus >= 0)}

\textbf{5º erro -- atribuição com \texttt{=} em vez de \texttt{:=}}

\begin{lstlisting}[language=Java]
class Teste4{
    // Outro programa de teste
    int idade;
    string nome, texto;
    float bonus, salario;
    {
    repeat{
        write("Digite o seu nome: ");
        read(nome);
        write("Digite o seu sobrenome");
        read(sobrenome);
        write("Digite a sua idade: ");
        read (idade);
        write("Digite o salario: ");
        read (salrio);
        write("Digite o bonus: ");
        read (bonus);
        if ((idade>0) and (salario > 0) and (bonus >= 0)){
            salarioLiquido = +salario salario*0.11) + (bonus
            109.8);
            texto = nome + " " + sobrenome + ":";
            write (text);
            write("Salario liquido: ");
            write(salario);
        };
    {until (idade <=0 or Salario <=0 or bonus <0);
    }
}
}
\end{lstlisting}

Código 9.4: Teste10\_corrigido3.txt

\begin{lstlisting}[language=bash]
> java -cp out Main compilador\Teste10_corrigido3.txt
Erro sintatico na linha 19: Esperado :=, encontrado '='
\end{lstlisting}

Análise: \texttt{salarioLiquido = ...} usa \texttt{=} em vez de \texttt{:=}.

Correção: \texttt{=} $\to$ \texttt{:=}

\textbf{6º erro -- expressão inválida após \texttt{:=}}

\begin{lstlisting}[language=Java]
class Teste4{
    // Outro programa de teste
    int idade;
    string nome, texto;
    float bonus, salario;
    {
    repeat{
        write("Digite o seu nome: ");
        read(nome);
        write("Digite o seu sobrenome");
        read(sobrenome);
        write("Digite a sua idade: ");
        read (idade);
        write("Digite o salario: ");
        read (salrio);
        write("Digite o bonus: ");
        read (bonus);
        if ((idade>0) and (salario > 0) and (bonus >= 0)){
            salarioLiquido := +salario salario*0.11) + (bonus
            109.8);
            texto = nome + " " + sobrenome + ":";
            write (text);
            write("Salario liquido: ");
            write(salario);
        };
    {until (idade <=0 or Salario <=0 or bonus <0);
    }
}
}
\end{lstlisting}

Código 9.5: Teste10\_corrigido4.txt

\begin{lstlisting}[language=bash]
> java -cp out Main compilador\Teste10_corrigido4.txt
Erro sintatico na linha 19: Fator esperado (identificador, constante, '(')
\end{lstlisting}

Análise: A expressão \texttt{+salario salario*0.11)} é inválida (operador unário \texttt{+} sem operando à esquerda, operandos justapostos). A linha foi removida na versão corrigida, simplificando o programa.

Correção: Linha \texttt{salarioLiquido := ...} removida; \texttt{texto =} $\to$ \texttt{texto :=}; \texttt{write(text)} $\to$ \texttt{write(texto)}

\textbf{7º erro -- sintaxe incorreta do repeat-until}

\begin{lstlisting}[language=Java]
class Teste4{
    // Outro programa de teste
    int idade;
    string nome, texto;
    float bonus, salario;
    {
    repeat{
        write("Digite o seu nome: ");
        read(nome);
        write("Digite o seu sobrenome");
        read(sobrenome);
        write("Digite a sua idade: ");
        read (idade);
        write("Digite o salario: ");
        read (salrio);
        write("Digite o bonus: ");
        read (bonus);
        if ((idade>0) and (salario > 0) and (bonus >= 0)){
            texto := nome + " " + sobrenome + ":";
            write (texto);
            write("Salario liquido: ");
            write(salario);
        };
    {until (idade <=0 or Salario <=0 or bonus <0);
    }
}
}
\end{lstlisting}

Código 9.6: Teste10\_corrigido5.txt

\begin{lstlisting}[language=bash]
> java -cp out Main compilador\Teste10_corrigido5.txt
Erro sintatico na linha 24: Esperado '}', encontrado '{'
\end{lstlisting}

Análise: O repeat-until deve ser \texttt{repeat \{ ... \} until (condição);}. O código usa \texttt{\{until} em vez de \texttt{\} until}.

Correção: \texttt{\{until} $\to$ \texttt{\} until}

\textbf{8º erro -- condição do until sem parênteses}

\begin{lstlisting}[language=Java]
class Teste4{
    // Outro programa de teste
    int idade;
    string nome, texto;
    float bonus, salario;
    {
    repeat{
        write("Digite o seu nome: ");
        read(nome);
        write("Digite o seu sobrenome");
        read(sobrenome);
        write("Digite a sua idade: ");
        read (idade);
        write("Digite o salario: ");
        read (salrio);
        write("Digite o bonus: ");
        read (bonus);
        if ((idade>0) and (salario > 0) and (bonus >= 0)){
            texto := nome + " " + sobrenome + ":";
            write (texto);
            write("Salario liquido: ");
            write(salario);
        };
    } until (idade <=0 or Salario <=0 or bonus <0);
    }
}
}
\end{lstlisting}

Código 9.7: Teste10\_corrigido6.txt

\begin{lstlisting}[language=bash]
> java -cp out Main compilador\Teste10_corrigido6.txt
Erro sintatico na linha 24: Esperado ')', encontrado '<='
\end{lstlisting}

Análise: A condição do \texttt{until} usa encadeamento de operadores relacionais com \texttt{or}. A gramática só aceita uma expressão relacional por \texttt{expression} (produção \texttt{expression' ::= relop\;simple-expr\;|\;$\lambda$}). O \texttt{or} é tratado como operador aditivo dentro de \texttt{simple-expr}, deixando o próximo \texttt{<=} fora da expressão, foi feita a correção de todos esses casos na linha.

Correção: \texttt{(idade <=0 or Salario <=0 or bonus <0)} $\to$ \texttt{((idade <= 0) or (salario <= 0) or (bonus < 0))}

\textbf{Teste aprovado}

\begin{lstlisting}[language=Java]
class Teste4{
    // Outro programa de teste
    int idade;
    string nome, texto;
    float bonus, salario;
    {
    repeat{
        write("Digite o seu nome: ");
        read(nome);
        write("Digite o seu sobrenome");
        read(sobrenome);
        write("Digite a sua idade: ");
        read (idade);
        write("Digite o salario: ");
        read (salario);
        write("Digite o bonus: ");
        read (bonus);
        if ((idade>0) and (salario > 0) and (bonus >= 0)){
            texto := nome + " " + sobrenome + ":";
            write (texto);
            write("Salario liquido: ");
            write(salario);
        };
    } until ((idade <= 0) or (salario <= 0) or (bonus < 0));
    }
}
\end{lstlisting}

Código 10: Teste10\_corrigido.txt

\begin{lstlisting}[language=bash]
> java -cp out Main compilador\Teste10_corrigido.txt
Programa sintaticamente correto.
\end{lstlisting}

\subsection{Teste 11 --- Fluxo de controle aninhado}

\textbf{1º erro -- atribuição dentro da declaração}

\begin{lstlisting}[language=Java]
class A {
    /** Verificando fluxo de controle
    Programa com if e while aninhados **/
    int i;
    float soma, altura, media, maior:=0, menor:=3;
    int soma;
    {
        soma := 0;

        write("Quantos dados deseja informar?");
        read(qtd);

        IF (qtd>=2){
            I:=0;
            do{
                do{
                    write("Altura: ");
                    read(altura);
                }while (altura <= 0 or altura >=2.5);

                soma := soma+altura;
                i := i + 1;
                if (altura > maior){
                    maior = altura;
                };
                if (altura < menor){
                    menor = Altura;
                };
            }while( i < qtd);

            media := soma / qtd;
            write("Media: " + media);

        }
        else{
            write("Quantidade invalida.");
        }
    }
}
\end{lstlisting}

Código 11: Teste11.txt -- com erros

\begin{lstlisting}[language=bash]
> java -cp out Main compilador\Teste11.txt
Erro sintatico na linha 5: Esperado ';', encontrado ':='
\end{lstlisting}

Análise: A declaração \texttt{float ... maior:=0, menor:=3;} contém \texttt{:=}, que não é permitido em declarações (apenas identificadores separados por vírgula).

Correção: \texttt{maior:=0, menor:=3} $\to$ \texttt{maior, menor} (atribuições movidas para o corpo)

\textbf{2º erro -- \texttt{IF} maiúsculo}

\begin{lstlisting}[language=Java]
class A {
    /** Verificando fluxo de controle
    Programa com if e while aninhados **/
    int i;
    float soma, altura, media, maior, menor;
    int soma;
    {
        soma := 0;
        maior := 0;
        menor := 3;

        write("Quantos dados deseja informar?");
        read(qtd);

        IF (qtd>=2){
            I:=0;
            do{
                do{
                    write("Altura: ");
                    read(altura);
                }while (altura <= 0 or altura >=2.5);

                soma := soma+altura;
                i := i + 1;
                if (altura > maior){
                    maior = altura;
                };
                if (altura < menor){
                    menor = Altura;
                };
            }while( i < qtd);

            media := soma / qtd;
            write("Media: " + media);

        }
        else{
            write("Quantidade invalida.");
        }
    }
}
\end{lstlisting}

Código 11.1: Teste11\_corrigido1.txt

\begin{lstlisting}[language=bash]
> java -cp out Main compilador\Teste11_corrigido1.txt
Erro sintatico na linha 15: Esperado :=, encontrado '('
\end{lstlisting}

Análise: \texttt{IF} (maiúsculo) é tratado como identificador (variável), não como palavra reservada \texttt{if}. O parser espera \texttt{:=} após um identificador.

Correção: \texttt{IF} $\to$ \texttt{if}

\textbf{3º erro -- condição do while sem parênteses}

\begin{lstlisting}[language=Java]
class A {
    /** Verificando fluxo de controle
    Programa com if e while aninhados **/
    int i;
    float soma, altura, media, maior, menor;
    int soma;
    {
        soma := 0;
        maior := 0;
        menor := 3;

        write("Quantos dados deseja informar?");
        read(qtd);

        if (qtd>=2){
            I:=0;
            do{
                do{
                    write("Altura: ");
                    read(altura);
                }while (altura <= 0 or altura >=2.5);

                soma := soma+altura;
                i := i + 1;
                if (altura > maior){
                    maior = altura;
                };
                if (altura < menor){
                    menor = Altura;
                };
            }while( i < qtd);

            media := soma / qtd;
            write("Media: " + media);

        }
        else{
            write("Quantidade invalida.");
        }
    }
}
\end{lstlisting}

Código 11.2: Teste11\_corrigido2.txt

\begin{lstlisting}[language=bash]
> java -cp out Main compilador\Teste11_corrigido2.txt
Erro sintatico na linha 21: Esperado ')', encontrado '>='
\end{lstlisting}

Análise: A condição \texttt{(altura <= 0 or altura >= 2.5)} precisa de cada expressão relacional entre parênteses para o \texttt{or} funcionar corretamente.

Correção: \texttt{(altura <= 0 or altura >= 2.5)} $\to$ \texttt{((altura <= 0) or (altura >= 2.5))}

\textbf{4º erro -- atribuição com \texttt{=} em vez de \texttt{:=}}

\begin{lstlisting}[language=Java]
class A {
    /** Verificando fluxo de controle
    Programa com if e while aninhados **/
    int i;
    float soma, altura, media, maior, menor;
    int soma;
    {
        soma := 0;
        maior := 0;
        menor := 3;

        write("Quantos dados deseja informar?");
        read(qtd);

        if (qtd>=2){
            I:=0;
            do{
                do{
                    write("Altura: ");
                    read(altura);
                }while ((altura <= 0) or (altura >= 2.5));

                soma := soma+altura;
                i := i + 1;
                if (altura > maior){
                    maior = altura;
                };
                if (altura < menor){
                    menor = Altura;
                };
            }while( i < qtd);

            media := soma / qtd;
            write("Media: " + media);

        }
        else{
            write("Quantidade invalida.");
        }
    }
}
\end{lstlisting}

Código 11.3: Teste11\_corrigido3.txt

\begin{lstlisting}[language=bash]
> java -cp out Main compilador\Teste11_corrigido3.txt
Erro sintatico na linha 26: Esperado :=, encontrado '='
\end{lstlisting}

Análise: \texttt{maior = altura;} usa \texttt{=} em vez de \texttt{:=}.

Correção: \texttt{maior = altura;} $\to$ \texttt{maior := altura;}

\textbf{5º erro -- atribuição com \texttt{=} (segunda ocorrência)}

\begin{lstlisting}[language=Java]
class A {
    /** Verificando fluxo de controle
    Programa com if e while aninhados **/
    int i;
    float soma, altura, media, maior, menor;
    int soma;
    {
        soma := 0;
        maior := 0;
        menor := 3;

        write("Quantos dados deseja informar?");
        read(qtd);

        if (qtd>=2){
            I:=0;
            do{
                do{
                    write("Altura: ");
                    read(altura);
                }while ((altura <= 0) or (altura >= 2.5));

                soma := soma+altura;
                i := i + 1;
                if (altura > maior){
                    maior := altura;
                };
                if (altura < menor){
                    menor = Altura;
                };
            }while( i < qtd);

            media := soma / qtd;
            write("Media: " + media);

        }
        else{
            write("Quantidade invalida.");
        }
    }
}
\end{lstlisting}

Código 11.4: Teste11\_corrigido4.txt

\begin{lstlisting}[language=bash]
> java -cp out Main compilador\Teste11_corrigido4.txt
Erro sintatico na linha 29: Esperado :=, encontrado '='
\end{lstlisting}

Análise: \texttt{menor = Altura;} usa \texttt{=} em vez de \texttt{:=}.

Correção: \texttt{menor = Altura;} $\to$ \texttt{menor := Altura;} (a correção de capitalização \texttt{Altura} $\to$ \texttt{altura} é semântica, não sintática)

\textbf{6º erro -- ponto e vírgula ausente após if-else}

\begin{lstlisting}[language=Java]
class A {
    /** Verificando fluxo de controle
    Programa com if e while aninhados **/
    int i;
    float soma, altura, media, maior, menor;
    int soma;
    {
        soma := 0;
        maior := 0;
        menor := 3;

        write("Quantos dados deseja informar?");
        read(qtd);

        if (qtd>=2){
            I:=0;
            do{
                do{
                    write("Altura: ");
                    read(altura);
                }while ((altura <= 0) or (altura >= 2.5));

                soma := soma+altura;
                i := i + 1;
                if (altura > maior){
                    maior := altura;
                };
                if (altura < menor){
                    menor := Altura;
                };
            }while( i < qtd);

            media := soma / qtd;
            write("Media: " + media);

        }
        else{
            write("Quantidade invalida.");
        }
    }
}
\end{lstlisting}

Código 11.5: Teste11\_corrigido5.txt

\begin{lstlisting}[language=bash]
> java -cp out Main compilador\Teste11_corrigido5.txt
Erro sintatico na linha 40: Esperado ';', encontrado '}'
\end{lstlisting}

Análise: O if-else é um comando na \texttt{stmt-list} e exige \texttt{;} ao final.

Correção: Adicionado \texttt{;} após o \texttt{\}} de fechamento do else

\textbf{Teste aprovado}

\begin{lstlisting}[language=Java]
class A {
    /** Verificando fluxo de controle
    Programa com if e while aninhados **/
    int i;
    float soma, altura, media, maior, menor;
    {
        soma := 0;
        maior := 0;
        menor := 3;

        write("Quantos dados deseja informar?");
        read(qtd);

        if (qtd>=2){
            i := 0;
            do{
                do{
                    write("Altura: ");
                    read(altura);
                }while ((altura <= 0) or (altura >= 2.5));

                soma := soma+altura;
                i := i + 1;
                if (altura > maior){
                    maior := altura;
                };
                if (altura < menor){
                    menor := altura;
                };
            }while( i < qtd);

            media := soma / qtd;
            write("Media: " + media);
        }
        else{
            write("Quantidade invalida.");
        };
    }
}
\end{lstlisting}

Código 12: Teste11\_corrigido.txt

\begin{lstlisting}[language=bash]
> java -cp out Main compilador\Teste11_corrigido.txt
Programa sintaticamente correto.
\end{lstlisting}

\subsection{Teste 12 --- If aninhado com repeat-until}

\textbf{1º erro -- corpo sem chaves de abertura e fechamento}

\begin{lstlisting}[language=Java]
class Teste6 {
    int numero, soma;
    string mensagem;

    soma := 0;

    repeat {
        write("Digite um numero (negativo para sair): ");
        read(numero);

        if (numero >= 0) {
            if (numero % 2 = 0) {
                mensagem := " par";
            } else {
                mensagem := " impar";
            }
            write("O numero " + numero + " eh" + mensagem);
            soma := soma + numero;
        }
    } until (numero < 0);

    write("Soma total: ");
    write(soma);
}
\end{lstlisting}

Código 13: Teste12.txt -- com erros

\begin{lstlisting}[language=bash]
> java -cp out Main compilador\Teste12.txt
Erro sintatico na linha 5: Esperado '{', encontrado 'soma'
\end{lstlisting}

Análise: Após \texttt{string mensagem;}, o corpo espera \texttt{\{} mas encontra \texttt{soma := 0;}.

Correção: Adicionado par de chaves \texttt{\{} \texttt{\}} (abertura e fechamento) para o corpo, antes e depois dos comandos

\textbf{2º erro -- ponto e vírgula ausente no if-else interno}

\begin{lstlisting}[language=Java]
class Teste6 {
    int numero, soma;
    string mensagem;
    {
    soma := 0;

    repeat {
        write("Digite um numero (negativo para sair): ");
        read(numero);

        if (numero >= 0) {
            if (numero % 2 = 0) {
                mensagem := " par";
            } else {
                mensagem := " impar";
            }
            write("O numero " + numero + " eh" + mensagem);
            soma := soma + numero;
        }
    } until (numero < 0);

    write("Soma total: ");
    write(soma);
    }
}
\end{lstlisting}

Código 13.1: Teste12\_corrigido1.txt

\begin{lstlisting}[language=bash]
> java -cp out Main compilador\Teste12_corrigido1.txt
Erro sintatico na linha 17: Esperado ';', encontrado 'write'
\end{lstlisting}

Análise: O if-else interno (\texttt{if (numero \% 2 = 0) \{ ... \} else \{ ... \}}) não possui \texttt{;} ao final.

Correção: Adicionado \texttt{;} após o \texttt{\}} do else interno

\textbf{3º erro -- ponto e vírgula ausente no if-else externo}

\begin{lstlisting}[language=Java]
class Teste6 {
    int numero, soma;
    string mensagem;
    {
    soma := 0;

    repeat {
        write("Digite um numero (negativo para sair): ");
        read(numero);

        if (numero >= 0) {
            if (numero % 2 = 0) {
                mensagem := " par";
            } else {
                mensagem := " impar";
            };
            write("O numero " + numero + " eh" + mensagem);
            soma := soma + numero;
        }
    } until (numero < 0);

    write("Soma total: ");
    write(soma);
    }
}
\end{lstlisting}

Código 13.2: Teste12\_corrigido2.txt

\begin{lstlisting}[language=bash]
> java -cp out Main compilador\Teste12_corrigido2.txt
Erro sintatico na linha 20: Esperado ';', encontrado '}'
\end{lstlisting}

Análise: O if-else externo também precisa de \texttt{;} para terminar o comando na \texttt{stmt-list}.

Correção: Adicionado \texttt{;} após o \texttt{\}} do if externo

\textbf{Teste aprovado}

\begin{lstlisting}[language=Java]
class Teste6 {
    int numero, soma;
    string mensagem;
    {
    soma := 0;

    repeat {
        write("Digite um numero (negativo para sair): ");
        read(numero);

        if (numero >= 0) {
            if (numero % 2 = 0) {
                mensagem := " par";
            } else {
                mensagem := " impar";
            };
            write("O numero " + numero + " eh" + mensagem);
            soma := soma + numero;
        };
    } until (numero < 0);

    write("Soma total: ");
    write(soma);
    }
}
\end{lstlisting}

Código 14: Teste12\_corrigido.txt

\begin{lstlisting}[language=bash]
> java -cp out Main compilador\Teste12_corrigido.txt
Programa sintaticamente correto.
\end{lstlisting}

\section{Conclusão}

A adequação da gramática para $LL(1)$ viabilizou o desenvolvimento de um analisador sintático descendente recursivo determinístico, com varredura linear e sem retrocesso. O uso dos conjuntos $FIRST$ e $FOLLOW$ guiou a implementação das rotinas de parsing, resultando em mensagens de erro precisas com linha e token causador. A organização em pacotes separados (\texttt{lexico}, \texttt{sintatico}, \texttt{tabela}) mantém a modularidade do código para as etapas futuras de análise semântica.

\end{document}
